\documentclass{article}

% ---------------------------------------------------------------------
% Preámbulo alineado al de teoricas/guia_tads/main.tex y guias/macros.tex
% (misma fuente, márgenes, encabezado con logo, título en versalitas,
% cajas de concepto/regla/invariante). Self-contained.
% ---------------------------------------------------------------------
\usepackage{amsmath}
\usepackage{amssymb}
\usepackage{graphicx}
\usepackage{minted}
\usepackage{clrscode3e}
\usepackage{booktabs}
\usepackage[spanish,es-noindentfirst,es-nosectiondot,es-nolists,
es-noshorthands,es-lcroman,es-tabla]{babel}
\usepackage{lmodern}
\usepackage[T1]{fontenc}
\usepackage[utf8]{inputenc}
\usepackage{textcomp}
\usepackage{microtype}
\usepackage[svgnames]{xcolor}
\usepackage[hidelinks,urlcolor=blue]{hyperref}
\hypersetup{colorlinks=true, allcolors=Navy, pdfstartview={XYZ null null 1}}
\usepackage[margin=2.5cm,headsep=0.8cm,headheight=1.3cm,footskip=1cm]{geometry}
\usepackage{tcolorbox}
\usepackage{enumitem}
\usepackage{fancyhdr}

\setlength{\parindent}{0pt}
\setlength{\parskip}{0.4\baselineskip}

\newcommand{\TODO}[1]{\textcolor{red}{[TODO: #1]}}

% Indicador de TAD para cada operación (igual que en la Guía de TADs).
\newcommand{\OP}[2]{\proc{#1}_{\text{#2}}}

\newminted{cpp}{autogobble,tabsize=4,fontsize=\small}
\newmintinline{cpp}{}

\newtcolorbox{concepto}[1]{
  colback=blue!5, colframe=blue!55!black, colbacktitle=blue!15,
  coltitle=black, fonttitle=\bfseries, boxrule=0.6pt, arc=2pt, title={#1},
}
\newtcolorbox{regla}[1]{
  colback=orange!5, colframe=orange!65!black, colbacktitle=orange!20,
  coltitle=black, fonttitle=\bfseries, boxrule=0.6pt, arc=2pt, title={#1},
}
\newtcolorbox{invariante}[1]{
  colback=green!5, colframe=green!40!black, colbacktitle=green!15,
  coltitle=black, fonttitle=\bfseries, boxrule=0.6pt, arc=2pt, title={#1},
}

\setlist[itemize]{topsep=2pt, itemsep=2pt}
\setlist[enumerate]{topsep=2pt, itemsep=2pt}

%%%% Encabezado y pie de página %%%%
\fancypagestyle{plain}{%
\fancyhf{}
\fancyfoot[C]{\thepage}
\fancyhead[L]{\raisebox{-0.15cm}{\includegraphics[height=0.8cm]{logoUdesa.png}}}
\fancyhead[R]{\footnotesize Algoritmos y Estructuras de Datos - IIA - UdeSA}
}
\pagestyle{plain}

\newcommand{\titulo}[2]{%
\begin{center}%
{\Large\textsc{Trabajo Práctico #1: #2}}
\end{center}}

\begin{document}

\titulo{3}{Goleadores del Mundial}

\begin{center}
\begin{minipage}{0.85\textwidth}
\small
\textbf{Tema:} tablas de hash y TAD Diccionario. \hfill
\textbf{Lenguaje:} C++. \\
\textbf{Modalidad:} en grupos de a 2. \hfill
\textbf{Fecha de entrega:} 5 de noviembre, 23:59.
\end{minipage}
\end{center}

% =======================================================================
\section{Objetivo}
% =======================================================================

En este trabajo vamos a representar un diccionario de jugadores y goles usando una \textbf{tabla de hash}, para que las operaciones de inserción, búsqueda y eliminación sean más rápidas que con un árbol balanceado (ver TAD Diccionario representado sobre un AVL en la teórica con
operaciones en $O(\log n)$).
Lo vamos a usar para llevar la tabla de
goleadores de un Mundial.

El trabajo se organiza en \textbf{tres TADs}, cada uno implementado
usando únicamente la interfaz del que tiene debajo:

\begin{center}
\begin{tabular}{@{}lll@{}}
\toprule
TAD & ¿Genérico? & Se representa con \\
\midrule
\texttt{Goleadores} & No: jugador (\texttt{string}) $\to$ goles (\texttt{int}) & un \texttt{Diccionario<string, int>} \\
\texttt{Diccionario<K, V>} & Sí & una \texttt{TablaHash<K, V>} \\
\texttt{TablaHash<K, V>} & Sí & un arreglo de \emph{buckets} \\
\bottomrule
\end{tabular}
\end{center}

Para cada TAD se presenta la \textbf{interfaz} (con su contrato y su costo) y
la \textbf{representación} (con su invariante). El objetivo es realizar la implementación de estos tipos de datos.

Además, \texttt{TablaHash} y \texttt{Diccionario} ofrecen cada uno un
\textbf{iterador} para recorrer su contenido (sección~\ref{sec:iterador}).

\begin{regla}{Sobre las complejidades de este enunciado}
En todas las cotas, $n$ representa la cantidad de claves guardadas en la
estructura en ese momento. Cuando dice \textbf{en promedio} se asume la hipótesis de \textbf{hashing uniforme} de la teórica: la función de hash reparte las claves de manera uniforme entre los buckets.
Cuando dice \textbf{amortizado} nos referimos al costo por operación a
lo largo de una secuencia de operaciones: alguna puede ser cara (la que
redimensiona), pero repartido entre todas el costo es el indicado. Es la
misma cuenta de la teórica para $\proc{InsertarAlFinal}$ del arreglo
dinámico y para la redimensión de la tabla de hash.
Salvo que se aclare, se supone $T_h \in O(1)$ (con $T_h$ el costo de
calcular la función de hash de una clave, como en la teórica) y que
comparar dos claves cuesta $O(1)$.

Lo marcado como \textbf{optativo} no es necesario para aprobar.
\end{regla}

% =======================================================================
\section{TAD Iterador}
\label{sec:iterador}
% =======================================================================

Es el TAD Iterador de la teórica: abstrae el \textbf{recorrido} de una
estructura, sin que quien lo usa conozca su representación interna. En
este trabajo lo que se recorre son pares $\langle clave, valor \rangle$,
así que $\proc{Actual}$ se separa en dos operaciones.

\subsection{Interfaz}

\begin{itemize}
\item $\OP{CrearIterador}{Iterador}(estructura)$: crea un iterador sobre $estructura$. \\
  {\footnotesize \textbf{Pre:} --- \quad \textbf{Post:} devuelve un iterador al primer par de $estructura$ (si no tiene pares, el iterador ya está al final).}
\item $\OP{HayMas}{Iterador}(it)$: indica si quedan pares por recorrer. \\
  {\footnotesize \textbf{Pre:} --- \quad \textbf{Post:} devuelve si quedan pares por recorrer, sin modificar $it$.}
\item $\OP{Clave}{Iterador}(it)$: la clave del par actual. \\
  {\footnotesize \textbf{Pre:} $\OP{HayMas}{Iterador}(it)$. \quad \textbf{Post:} devuelve la clave del par actual, sin modificar $it$.}
\item $\OP{Valor}{Iterador}(it)$: el valor del par actual. \\
  {\footnotesize \textbf{Pre:} $\OP{HayMas}{Iterador}(it)$. \quad \textbf{Post:} devuelve el valor del par actual, sin modificar $it$.}
\item $\OP{Avanzar}{Iterador}(it)$: pasa al par siguiente. \\
  {\footnotesize \textbf{Pre:} $\OP{HayMas}{Iterador}(it)$. \quad \textbf{Post:} $it$ pasa al par siguiente (o al final, si no hay más).}
\end{itemize}

Un recorrido completo pasa por \textbf{cada par exactamente una vez}, en
\textbf{algún orden} (no se puede asumir cuál). El iterador deja de ser
válido si se inserta o se elimina en la estructura mientras se lo usa.

El costo de cada operación depende de la estructura que se recorre, y
se indica en la sección de cada una.

\subsection{En C++}

No hay una única clase \texttt{Iterador}: cada estructura define la
suya, \textbf{anidada} dentro de su propia clase y con estas mismas
operaciones. El código que recorre es siempre igual; sólo cambia el tipo
del iterador:

\begin{cppcode}
TablaHash<std::string, int>::Iterador it = tabla.crear_iterador();
while (it.hay_mas()) {
    std::cout << it.clave() << ": " << it.valor() << std::endl;
    it.avanzar();
}
\end{cppcode}

En el esqueleto, la declaración de cada \texttt{Iterador} ya está
escrita: hay que completar su representación (en \texttt{TablaHash}) y
el cuerpo de sus operaciones.

% =======================================================================
\section{TAD TablaHash}
% =======================================================================

Una tabla de hash con \textbf{direccionamiento cerrado}: es un arreglo de
\emph{buckets}, donde cada bucket guarda todos los pares
$\langle clave, valor \rangle$ cuyo hash cae en esa posición.

Es la misma estructura de la teórica, pero \textbf{no la misma
interfaz}: allá la tabla era directamente el diccionario
($\proc{Insertar}$ reemplazaba el valor si la clave ya estaba, y
$\proc{Pertenece}$ y $\proc{Obtener}$ eran operaciones de la tabla). Acá la
tabla es la capa de abajo, con operaciones más simples y precondiciones
más fuertes; el comportamiento de diccionario lo agrega el TAD
Diccionario, usándola.

% Es la estructura ``de abajo'': no sabe qué es un diccionario. Por eso
% sus operaciones son simples y tienen \textbf{precondiciones fuertes}
% (quien la usa es responsable de cumplirlas).

\subsection{Interfaz}

Con $n$ la cantidad de claves guardadas y $c$ la capacidad (cantidad de
buckets) de la tabla:

\begin{itemize}
\item $\OP{TablaVacia}{TablaHash}(c_0)$: crea una tabla vacía. \hfill $O(c_0)$ \\
  {\footnotesize \textbf{Pre:} $c_0 \geq 1$. \quad \textbf{Post:} devuelve una tabla sin claves, con $c_0$ buckets.}
\item $\OP{Insertar}{TablaHash}(T, clave, valor)$: agrega un par nuevo. \hfill $O(1)$ amortizado, $O(n)$ en peor caso \\
  {\footnotesize \textbf{Pre:} $clave \notin T$. \quad \textbf{Post:} $T$ tiene además el par $\langle clave, valor \rangle$.}
\item $\OP{Eliminar}{TablaHash}(T, clave)$: quita $clave$ y su valor. \hfill $O(1)$ en promedio, $O(n)$ en peor caso \\
  {\footnotesize \textbf{Pre:} $clave \in T$. \quad \textbf{Post:} $T$ ya no tiene a $clave$.}
\item $\OP{Buscar}{TablaHash}(T, clave)$: busca el valor de $clave$. \hfill $O(1)$ en promedio, $O(n)$ en peor caso \\
  {\footnotesize \textbf{Pre:} --- \quad \textbf{Post:} si $clave \in T$, devuelve una referencia a su valor (se puede leer y modificar); si no, $\const{nil}$.}
\item $\OP{Cantidad}{TablaHash}(T)$: cantidad de claves guardadas. \hfill $O(1)$ \\
  {\footnotesize \textbf{Pre:} --- \quad \textbf{Post:} devuelve $n$.}
\item $\OP{Capacidad}{TablaHash}(T)$: cantidad de buckets. \hfill $O(1)$ \\
  {\footnotesize \textbf{Pre:} --- \quad \textbf{Post:} devuelve $c$.}
\item $\OP{Colisiones}{TablaHash}(T)$: cantidad de colisiones. \hfill $O(c)$ en peor caso \\
  {\footnotesize \textbf{Pre:} --- \quad \textbf{Post:} devuelve $n$ menos la cantidad de buckets no vacíos (es decir, cuántas claves comparten bucket con otra que llegó antes).}
\item $\OP{CrearIterador}{TablaHash}(T)$: crea un iterador sobre los pares de $T$. \hfill $O(c)$ en peor caso \\
  {\footnotesize \textbf{Pre:} --- \quad \textbf{Post:} devuelve un iterador al primer par de $T$ (ver sección~\ref{sec:iterador}).}
\end{itemize}

En la implementación, además, la tabla tiene que \textbf{liberar toda su memoria} al destruirse (destructor), en $O(n + c)$ en peor caso.

\subsection{Iterador de la tabla}

Las operaciones son las del TAD Iterador (sección~\ref{sec:iterador}),
con estos costos:

\begin{center}
\begin{tabular}{@{}ll@{}}
\toprule
Operación & Costo \\
\midrule
$\proc{HayMas}$, $\proc{Clave}$, $\proc{Valor}$ & $O(1)$ \\
$\proc{Avanzar}$ & $O(c)$ en peor caso \\
Recorrido completo ($\proc{CrearIterador}$ y $\proc{Avanzar}$ hasta el final) & $O(n + c)$ en peor caso \\
\bottomrule
\end{tabular}
\end{center}

A diferencia del iterador de la lista enlazada, acá $\proc{Avanzar}$
\textbf{no} es $O(1)$: puede tener que pasar por muchos buckets vacíos
hasta encontrar el par siguiente. Pero cada bucket y cada par se miran
una sola vez en todo el recorrido, y por eso el total es $O(n + c)$.

La representación del iterador también \textbf{la eligen ustedes} (¿qué
necesita recordar para saber dónde está y cómo seguir?), y también
tienen que escribir su invariante como comentario.

\subsection{Representación interna (e invariante)}

La representación \textbf{la eligen ustedes}: en particular, qué
estructura usan para cada bucket (una lista enlazada, un arreglo
dinámico, \ldots). Tienen que escribir su \textbf{invariante de
representación} como comentario en \texttt{tabla\_hash.h}. Cualquiera sea
la representación, tiene que cumplir:

\begin{invariante}{Lo que tiene que valer en cualquier representación}
\begin{itemize}
\item $c \geq 1$.
\item El par de cada clave está en el bucket número
  $h(clave) = \texttt{funcion\_hash}(clave) \bmod c$.
\item Ninguna clave aparece más de una vez.
\item $n \leq \texttt{FACTOR\_DE\_CARGA\_MAXIMO} \cdot c$ (es el $\alpha_{max}$ de la teórica).
\end{itemize}
\end{invariante}

El último punto es el que garantiza los costos: como el \textbf{factor
de carga} $\alpha = n / c$ está acotado por una constante, cada bucket
tiene $O(1)$ pares en promedio. Para mantenerlo, cuando una inserción lo
rompería hay que \textbf{redimensionar}: pasar a un arreglo de buckets
más grande (por ejemplo, del doble de capacidad) y reubicar todos los
pares (¡el bucket de cada clave cambia, porque cambia $c$!). No se pide
achicar la tabla al eliminar.

En el esqueleto, $\texttt{FACTOR\_DE\_CARGA\_MAXIMO} = 2$.

\subsection{Función de hash}

La tabla llama a \cppinline{funcion_hash(clave)}, que devuelve un
\cppinline{size_t}. Hay que implementar dos versiones (sobrecarga): una
para \cppinline{std::string}, en $O(m)$ con $m$ la longitud del string, y
otra para \cppinline{int}, en $O(1)$. Lo único obligatorio es que dos
claves iguales tengan siempre el mismo hash: \textbf{cuál} función usar
es parte del trabajo.

\subsection{Exploración: ¿qué función de hash conviene?}

No hay una única respuesta, así que hay que \textbf{probar}. El
programa \texttt{explorar\_hash.cpp} carga en una \texttt{TablaHash}
los jugadores del archivo de datos del grupo y muestra la cantidad de
claves, la capacidad final y las colisiones.

\begin{enumerate}
\item Proponer \textbf{al menos tres} funciones de hash distintas para
  strings (alguna deliberadamente mala, por ejemplo una que sólo mire el
  largo o la primera letra).
\item Para cada una, medir las colisiones sobre el archivo del grupo,
  con al menos dos capacidades iniciales distintas.
\item Volcar los resultados en una tabla y explicar en pocas líneas qué
  función eligieron y por qué. ¿Qué tienen en común las claves que
  colisionan con las funciones malas? ¿Cambia algo si la capacidad es
  potencia de 2?
\end{enumerate}

Como referencia: si la función repartiera las claves de manera
perfectamente uniforme, con $n$ claves y $c$ buckets se esperan
aproximadamente $n - c\,(1 - (1 - 1/c)^n)$ colisiones.

% =======================================================================
\section{TAD Diccionario}
% =======================================================================

La misma interfaz de la teórica, más la cantidad de claves. A
diferencia de la tabla, acá $\proc{Insertar}$ \textbf{no tiene
precondición}: si la clave ya estaba, reemplaza su valor.

\subsection{Interfaz}

\begin{itemize}
\item $\OP{DiccionarioVacio}{Diccionario}()$: crea un diccionario vacío. \hfill $O(1)$ \\
  {\footnotesize \textbf{Pre:} --- \quad \textbf{Post:} devuelve el diccionario vacío.}
\item $\OP{Insertar}{Diccionario}(D, clave, valor)$ \hfill $O(1)$ amortizado en promedio, $O(n)$ en peor caso \\
  {\footnotesize \textbf{Pre:} --- \quad \textbf{Post:} $D$ tiene $clave \to valor$ (reemplazando el valor anterior, si $clave$ ya estaba).}
\item $\OP{Pertenece}{Diccionario}(D, clave)$: ¿está $clave$? \hfill $O(1)$ en promedio, $O(n)$ en peor caso \\
  {\footnotesize \textbf{Pre:} --- \quad \textbf{Post:} devuelve si $clave$ tiene un valor asociado en $D$.}
\item $\OP{Obtener}{Diccionario}(D, clave)$: el valor asociado a $clave$. \hfill $O(1)$ en promedio, $O(n)$ en peor caso \\
  {\footnotesize \textbf{Pre:} $\OP{Pertenece}{Diccionario}(D, clave)$. \quad \textbf{Post:} devuelve el valor asociado a $clave$.}
\item $\OP{Eliminar}{Diccionario}(D, clave)$: quita $clave$ (y su valor). \hfill $O(1)$ en promedio, $O(n)$ en peor caso \\
  {\footnotesize \textbf{Pre:} $\OP{Pertenece}{Diccionario}(D, clave)$. \quad \textbf{Post:} $D$ ya no tiene a $clave$.}
\item $\OP{Cantidad}{Diccionario}(D)$: cantidad de claves. \hfill $O(1)$ \\
  {\footnotesize \textbf{Pre:} --- \quad \textbf{Post:} devuelve $n$.}
\item $\OP{CrearIterador}{Diccionario}(D)$: crea un iterador sobre las asociaciones de $D$. \hfill $O(c)$ en peor caso \\
  {\footnotesize \textbf{Pre:} --- \quad \textbf{Post:} devuelve un iterador a la primera asociación de $D$ (ver sección~\ref{sec:iterador}).}
\end{itemize}

En $\OP{CrearIterador}{Diccionario}$, $c$ es la capacidad de la tabla de
hash subyacente. Las operaciones del iterador del diccionario tienen los
mismos costos que las del iterador de la tabla.

En C++: \cppinline{bool pertenece(const K& clave)} y
\cppinline{V obtener(const K& clave)}.

\subsection{Representación interna (e invariante)}

\begin{invariante}{Invariante de representación}
$D.tabla$ es una \texttt{TablaHash} (que cumple su propio invariante)
cuyos pares $\langle clave, valor \rangle$ son exactamente las
asociaciones de $D$.
\end{invariante}

El diccionario \textbf{sólo} puede usar la interfaz pública de
\texttt{TablaHash}: no conoce los buckets. Presten atención a que cada
llamada a la tabla cumpla su precondición.

Lo mismo vale para su iterador, que se representa con un iterador de la
tabla:

\begin{invariante}{Invariante de representación del iterador}
$it.de\_la\_tabla$ es un iterador de $D.tabla$, y su par actual es la
asociación actual de $it$.
\end{invariante}

% =======================================================================
\section{TAD Goleadores}
% =======================================================================

La tabla de goleadores de un torneo: para cada jugador que convirtió,
cuántos goles lleva. No es genérico.

\subsection{Interfaz}

Con $n$ la cantidad de jugadores distintos y $m$ la longitud del nombre
del jugador (acá las claves son strings, así que $T_h \in O(m)$ y
comparar dos claves también cuesta $O(m)$):

\begin{itemize}
\item $\OP{SinGoles}{Goleadores}()$: crea una tabla de goleadores vacía. \hfill $O(1)$ \\
  {\footnotesize \textbf{Pre:} --- \quad \textbf{Post:} ningún jugador tiene goles.}
\item $\OP{AnotarGol}{Goleadores}(G, jugador)$: le suma un gol a $jugador$. \hfill $O(m)$ amortizado en promedio \\
  {\footnotesize \textbf{Pre:} --- \quad \textbf{Post:} $jugador$ tiene un gol más que antes (si no tenía ninguno, pasa a tener $1$).}
\item $\OP{Goles}{Goleadores}(G, jugador)$: cuántos goles lleva $jugador$. \hfill $O(m)$ en promedio \\
  {\footnotesize \textbf{Pre:} --- \quad \textbf{Post:} devuelve los goles de $jugador$ ($0$ si no anotó ninguno). No modifica $G$.}
\item $\OP{CantidadDeJugadores}{Goleadores}(G)$: cuántos jugadores anotaron. \hfill $O(1)$ \\
  {\footnotesize \textbf{Pre:} --- \quad \textbf{Post:} devuelve la cantidad de jugadores con al menos un gol.}
\item $\OP{MaximoGoleador}{Goleadores}(G)$: quién hizo más goles. \hfill $O(m)$ en peor caso (no depende de $n$) \\
  {\footnotesize \textbf{Pre:} $\OP{CantidadDeJugadores}{Goleadores}(G) > 0$. \quad \textbf{Post:} devuelve el jugador con más goles; si hay varios empatados, el que llegó primero a esa cantidad.}
\item \textbf{(Optativo)} $\OP{ImprimirPorNombre}{Goleadores}(G)$ \hfill $O(n \log n)$ comparaciones, en peor caso \\
  {\footnotesize \textbf{Pre:} --- \quad \textbf{Post:} imprime una línea ``\texttt{jugador: goles}'' por jugador, en orden alfabético.}
\item \textbf{(Optativo)} $\OP{ImprimirPorGoles}{Goleadores}(G)$ \hfill $O(n \log n)$ comparaciones, en peor caso \\
  {\footnotesize \textbf{Pre:} --- \quad \textbf{Post:} imprime una línea ``\texttt{jugador: goles}'' por jugador, de más a menos goles; a igual cantidad de goles, en orden alfabético.}
\end{itemize}

El orden alfabético es el del operador \cppinline{<} entre
\cppinline{std::string}. Los nombres de los archivos de datos no
tienen tildes ni eñes.

\subsection{Representación interna (e invariante)}

$G = \langle G.goles\_por\_jugador,\; G.maximo \rangle$, donde
$G.maximo$ es \textbf{únicamente un string} (no se guardan aparte sus
goles).

\begin{invariante}{Invariante de representación}
\begin{itemize}
\item $G.goles\_por\_jugador$ es un \texttt{Diccionario<string, int>}
  cuyas claves son exactamente los jugadores con al menos un gol, y el
  valor de cada uno es su cantidad de goles (siempre $\geq 1$).
\item Si hay al menos un jugador, $G.maximo$ es el nombre de un jugador
  con la mayor cantidad de goles.
\end{itemize}
\end{invariante}

Igual que con $C.maximo$ en el Conjunto sobre ABB, el máximo no se
recalcula: \textbf{cada operación} que modifica $G$ tiene que dejarlo
consistente con el diccionario.

Para la parte optativa: una tabla de hash \textbf{no guarda las claves
en ningún orden útil}, así que para imprimir los listados hay que volcar
el contenido del diccionario a un arreglo (recorriéndolo con su
iterador) y \textbf{ordenarlo}, con un algoritmo de ordenamiento
implementado por ustedes.

% =======================================================================
\section{El programa}
% =======================================================================

Cada grupo recibe un archivo de texto con los goles de un Mundial: una
línea por gol, con el nombre del jugador que lo hizo. El programa
(\texttt{main.cpp}) lee esa lista de la \textbf{entrada estándar}; la
lectura ya está resuelta en el esqueleto, hasta dejar cada nombre en una
variable. Falta crear el \texttt{Goleadores}, cargarle los goles e
imprimir el resumen.

Por ejemplo, para la final de 2022 (\texttt{datos/ejemplo.txt}):

\begin{minipage}[t]{0.3\textwidth}
\textbf{Entrada}
\begin{verbatim}
Messi
Di Maria
Mbappe
Mbappe
Messi
Mbappe
\end{verbatim}
\end{minipage}
\hfill
\begin{minipage}[t]{0.62\textwidth}
\textbf{Salida}
\begin{verbatim}
Jugadores con goles: 3
Maximo goleador: Mbappe (3 goles)
\end{verbatim}
\textbf{Salida adicional de la parte optativa}
\begin{verbatim}

Goleadores por nombre:
Di Maria: 1
Mbappe: 3
Messi: 2

Goleadores por goles:
Mbappe: 3
Messi: 2
Di Maria: 1
\end{verbatim}
\end{minipage}

La salida tiene que respetar \textbf{exactamente} ese formato. Se
corre así:

\begin{verbatim}
./goleadores < datos/ejemplo.txt
\end{verbatim}

% =======================================================================
\section{Archivos y reglas}
% =======================================================================

\begin{center}
\begin{tabular}{@{}ll@{}}
\toprule
Archivo & Qué hay que hacer \\
\midrule
\texttt{funcion\_hash.h} / \texttt{.cpp} & Implementar las dos funciones de hash. \\
\texttt{tabla\_hash.h} & Elegir la representación e implementar \texttt{TablaHash<K, V>} y su iterador. \\
\texttt{diccionario.h} & Implementar las operaciones de \texttt{Diccionario<K, V>} y de su iterador. \\
\texttt{goleadores.h} / \texttt{.cpp} & Implementar las operaciones de \texttt{Goleadores}. \\
\texttt{main.cpp} & Completar el programa (la lectura ya está). \\
\texttt{explorar\_hash.cpp} & Completar el programa de exploración (la lectura ya está). \\
\texttt{tests.cpp} & Pruebas de la parte obligatoria. Agreguen las suyas. \\
\texttt{Makefile} & \texttt{make}, \texttt{make test}, \texttt{make ejemplo}, \texttt{make valgrind}; \\
 & \texttt{make ejemplo-optativo} (los listados de la parte optativa). \\
\bottomrule
\end{tabular}
\end{center}

\texttt{TablaHash} y \texttt{Diccionario} son clases template, como la
lista enlazada del laboratorio: su implementación va en el mismo
\texttt{.h}. Las operaciones de cada \texttt{Iterador} se completan
directamente dentro de la clase anidada. Lo único que se puede hacer con
una clave es copiarla, compararla con \cppinline{==} y pasársela a
\cppinline{funcion_hash}; con un valor, copiarlo.

\begin{regla}{Reglas}
\begin{itemize}
\item \textbf{No se puede cambiar la interfaz pública} de ninguna de las
  clases. La parte privada de \texttt{TablaHash} y de su iterador es
  libre; la de \texttt{Diccionario}, su iterador y \texttt{Goleadores}
  es la de este enunciado
  (se pueden agregar funciones auxiliares privadas).
\item De la biblioteca estándar sólo se puede usar
  \texttt{<string>}, \texttt{<iostream>}, \texttt{<cstddef>} y
  \texttt{<vector>}. En particular, \textbf{no} se puede usar
  \texttt{std::map}, \texttt{std::unordered\_map}, \texttt{std::hash} ni
  \texttt{std::sort}.
\item Cada TAD usa sólo la interfaz pública del de abajo.
\item El programa no puede perder memoria ni acceder a memoria
  inválida (\texttt{make valgrind} tiene que terminar sin errores).
\item Tienen que cumplirse las complejidades de este enunciado.
\end{itemize}
\end{regla}

% =======================================================================
\section{Entrega}
% =======================================================================

Se entrega por el campus, hasta el 5 de noviembre a las 23:59, un único
archivo comprimido por grupo con:

\begin{itemize}
\item Todos los archivos del esqueleto, completos.
\item El archivo de datos que le tocó al grupo, como
  \texttt{datos/goles.txt}.
\item La salida del programa para ese archivo, como
  \texttt{datos/salida.txt}
  (\texttt{./goleadores < datos/goles.txt > datos/salida.txt}).
\item La exploración de funciones de hash, como
  \texttt{exploracion.txt}: la tabla de resultados y la explicación.
\end{itemize}

Para aprobar, el trabajo tiene que compilar sin errores, pasar las
pruebas de la cátedra (que incluyen a las de \texttt{tests.cpp}), no
tener errores de memoria, respetar las reglas de la sección anterior e
incluir la exploración.

\end{document}
